\subsection{An actual fractional mixed trace model}
The marked-trace gate in \fid{F54} has a genuine
one-place exception when arbitrary mixed radii are allowed.
The following construction retains the actual right module
and two specified trace laws. Its failure of stronger
comparisons can be checked on explicit operators.

\begin{finding}[\fid{F55}, Part 1: a regular one-place model]
\label{f:fractionalescape}
There is one finite localization of the algebraic cover
$v=t^2$ over $\Q(\sqrt2)$, containing $t=-1$ and all points
$t_n=2^{n/2}$, on which the actual rank-$64$ tensor module
has an invertible regular operator $F$ commuting with its
actual right action. With base $t_1=\sqrt2$, it satisfies
for every $n\ge1$
\begin{equation}\label{eq:fractionaltwotraces}
 \Tr F(t_n)=\Tr(F(\sqrt2)^n),\qquad
 \Tr\bigl(X(2^n)F(t_n)\bigr)
       =\Tr\bigl(X(2)F(\sqrt2)^n\bigr),
\end{equation}
where $X$ is the actual marked central element
\eqref{eq:actualmarkedX}.
The family is regular at the chosen collision $t=-1$,
where $v=1$, and retains its order-two collision.
It does not commute with all named ordinary left
operators, does not have full spectral retuning, and
fails the fixed $P_1$ and $P_1^2$ trace comparisons at
every $n>1$. Its operator entries have unavoidable poles
at the other place $t=+1$ above $v=1$.
\end{finding}

\begin{proof}
\par\smallskip\noindent\textbf{The actual summands.}\quad
Use the regular \fid{F21}/\fid{F43} bimodule isomorphism
\begin{align}
 M={}&(W_3\otimes W_3\oplus W_1\otimes W_1)
                          \otimes\epsilon_+
       \ \oplus\ (W_1\otimes W_1)\otimes\epsilon_-
       \notag\\
 &\oplus(W_3\otimes W_1\oplus W_1\otimes W_3)
                          \otimes N(\delta)
       \ \oplus\ (W_1\otimes W_1)\otimes N(D-2).
 \label{eq:fractionalactualmodule}
\end{align}
The right-isotypic multiplicities are $20,4,16,4$,
and their actual dimensions are $20,4,32,8$.
In the final summand the multiplicity space is exactly
$W_1^V\otimes W_1^W$.
Here $W_1$ denotes the $\mathfrak{sl}_2$-highest-weight-one
factor of the polynomial $\GL_2$ representation of
partition $(2,1)$, with its determinant twist retained.

Put
\[
 \delta=v+v^{-1},\qquad D=\delta^2,\qquad
 \Delta=\frac{(D-1)(D-4)}{D^2}.
\]
The element $X$ acts by $\Delta$ on the last rank-eight
summand and by zero on the others; on the actual regular
tensor lattice, $X(1)=0$.
The preceding bimodule comparison is defined over
\[
 S_0=\Q[v,v^{-1},\delta^{-1},(D-1)^{-1},(D-2)^{-1}].
\]
We pull back this actual comparison.
The rank-eight summand is selected by its regular
commutant decomposition, without dividing $X$ by
$\Delta$ at the collision.

\par\smallskip\noindent\textbf{One coefficient cover.}\quad
Set
\[
 s=\sqrt2>0,\qquad K_0=\Q(s),\qquad v=t^2,\qquad
 \eta=\Delta(2)=\frac{189}{625}.
\]
The base is $t=s$, the extension points are $t_n=s^n$,
and the chosen collision is $t=-1$.
Define
\begin{equation}\label{eq:fractionalfunctions}
 g(t)=2(1+t)^2,\qquad
 h(t)=\frac{\eta g(t)}{\Delta(t^2)},\qquad
 \lambda(t)=\frac{h(t)}2-1-2s,\qquad
 u(t)=1+\frac{g(t)-h(t)}{20}.
\end{equation}
The cancellation at $t=-1$ is exact:
\[
 \Delta(t^2)=
 \frac{(t-1)^2(t+1)^2(t^2+1)^2(t^8+t^4+1)}
      {(t^4+1)^4},
\]
\begin{equation}\label{eq:fractionalhcancellation}
 h(t)=\frac{2\eta(t^4+1)^4}
 {(t-1)^2(t^2+1)^2(t^8+t^4+1)}.
\end{equation}
In particular,
\[
\begin{array}{c|ccc}
 t&h(t)&u(t)&\lambda(t)\\ \hline
 s&6+4s&1&2\\
 -1&126/625&6187/6250&-562/625-2s .
\end{array}
\]
At the remaining small extension point,
\[
 h(2)=\frac{1503378}{203125},\qquad
 \lambda(2)=\frac{548564}{203125}-2s\ne0.
\]
The latter is nonzero because its first term is rational
and $s$ is irrational.

Let $R_0$ be the pullback of $S_0$ to $K_0[t,t^{-1}]$,
further localized at $t-1$ and $t^2+1$, and put
\begin{equation}\label{eq:fractionalbase}
 R=R_0[u^{-1},\lambda^{-1}].
\end{equation}
The expression for $h$ belongs to $R_0$, since
\[
 D(t^2)-1=\frac{t^8+t^4+1}{t^4}.
\]
This is one fixed finite localization of the coefficient
cover. Every inherited factor is nonzero at $v=1$
and at real $v>1$. The two new denominator factors
are nonzero at $t=-1$ and at every $t_n>1$.
The displayed values settle $u,\lambda$ at $t=-1$
and $\lambda$ at $n=1,2$.

For $n\ge3$, $t_n\ge2s$ and $0<\Delta(2^n)<1$, so
\[
 \lambda(t_n)>
 \eta(1+2s)^2-1-2s
 =\frac{1076-494s}{625}>0,
\]
where $s<2$ already gives a numerator greater than $88$.
To check $u$ at all $n>1$, $D(v)$ is strictly increasing
for $v>1$, and
\[
 \frac{d}{dD}\left(1-\frac5D+\frac4{D^2}\right)
      =\frac{5D-8}{D^3}>0\qquad(D>4).
\]
Thus $\Delta(2^n)>\eta$, whence
\begin{equation}\label{eq:fractionalmonotonicity}
 h(t_n)<g(t_n),\qquad u(t_n)>1\quad(n>1).
\end{equation}
All required points survive \eqref{eq:fractionalbase}.
At $t=-1$, $v-1=(t-1)(t+1)$ has a simple zero,
so $\Delta(t^2)$ still has a zero of order exactly two.
The actual coherent center and tensor module are
their original pullbacks at this place.

\par\smallskip\noindent\textbf{The actual automorphism.}\quad
On the entire rank-$20$ $\epsilon_+$ block in
\eqref{eq:fractionalactualmodule}, let $F=uI_{20}$.
This includes both the $W_3\otimes W_3$ and
$W_1\otimes W_1$ multiplicity summands of that scalar
right type.
On $\epsilon_-^4\oplus N(\delta)^{16}$, of total
dimension $36$, let $F$ be the identity.
On the final block, in the ordered multiplicity basis
\[
 g_+\otimes g_+,\quad g_+\otimes g_-,\quad
 g_-\otimes g_+,\quad g_-\otimes g_-,
\]
put
\begin{equation}\label{eq:fractionaloperator}
 F=\operatorname{diag}(1,s,s,\lambda(t))\otimes I_2.
\end{equation}
Every entry is regular over $R$, and
\[
 \det F=u^{20}(s^2\lambda)^2=4u^{20}\lambda^2
          \in R^\times.
\]
If $\iota$ is the regular $A$-linear isomorphism
from the original tensor module to
\eqref{eq:fractionalactualmodule}, the operator on the
original tensor module is explicitly
$\iota^{-1}F_{\rm block}\iota$, with the rightmost map
applied first.
Thus the comparison retains the actual named right
generators; it does not substitute a trace-equivalent
module. On the nonscalar blocks $F$ acts on the
multiplicity factor and the right action on $N(a)$.
It therefore commutes with that right action.
No division by $\Delta$ enters its definition.

\par\smallskip\noindent\textbf{The two trace identities.}\quad
On the final rank-eight sector,
\[
 \Tr(F|_{M_2})=2(1+s+s+\lambda)=h.
\]
Consequently
\begin{equation}\label{eq:fractionalalgebraictraces}
 \Tr_M(X(t^2)F(t))=\Delta(t^2)h(t)=\eta g(t),\qquad
 \Tr_M F(t)=20u(t)+36+h(t)=56+g(t).
\end{equation}
At $t=s$ the full base spectrum is
\[
 1^{\,58},\qquad s^{\,4},\qquad 2^{\,2}.
\]
On the marked sector the respective multiplicities
are $2,4,2$, and $X$ has eigenvalue $\eta$.
As $s^2=2$,
\[
 \Tr(F(s)^n)=58+4s^n+2\cdot2^n=56+2(1+s^n)^2,
\]
\[
 \Tr(X(2)F(s)^n)
       =\eta(2+4s^n+2\cdot2^n)=\eta\,2(1+s^n)^2.
\]
Substitution in \eqref{eq:fractionalalgebraictraces}
proves \eqref{eq:fractionaltwotraces} for every $n$.

\par\smallskip\noindent\textbf{Ordinary left action.}\quad
The actual \fid{F43} operator has
$E_V(g_-\otimes g_+)=g_+\otimes g_+$ on the final
multiplicity space. Hence
\[
 (F(s)E_V-E_VF(s))(g_-\otimes g_+)
       =(1-s)(g_+\otimes g_+)\ne0.
\]
Tensor with any nonzero right-factor vector.
This uses composition of actual operators on vectors;
the source's row-matrix/opposite-composition convention
is unchanged.
Thus ordinary left invariance fails explicitly, as
required by the both-actions gate in \fid{F54}.

\par\smallskip\noindent\textbf{Fixed generator words and spectra.}\quad
On $N(a)$, $\Tr P_1=D$.
On $\epsilon_+$ it is $D$ and on $\epsilon_-$ it is zero.
Therefore
\[
 \Tr(P_1(v)F(t))
 =D(v)\left(20u(t)+16+\frac{h(t)}2\right)
 =D(v)\left(36+g(t)-\frac{h(t)}2\right).
\]
At the base,
\[
 \Tr(P_1(2)F(s)^n)
  =D(2)(37+2s^n+2^n)
  =D(2)\left(36+\frac{g(t_n)}2\right).
\]
By \eqref{eq:fractionalmonotonicity} and
$D(2^n)>D(2)$, the first expression at $t_n$
is strictly larger for every $n>1$.
Multiplying by the respective $D$ factors, using
$P_1^2=DP_1$, gives the same strict failure for $P_1^2$.

Full spectral retuning also fails without retaining
any sector labels. The operator $F(s)^n$ has eigenvalue
$1$ with multiplicity exactly $58$.
For $n>1$, $F(t_n)$ has that eigenvalue with multiplicity
at most $40$: the $36$ unchanged dimensions, the two
fixed marked dimensions, and at most two additional
dimensions if $\lambda(t_n)=1$.
Here $u(t_n)>1$ and $s>1$ exclude all other sources.

\par\smallskip\noindent\textbf{The other place above one.}\quad
At $t=+1$, \eqref{eq:fractionalhcancellation} has a pole
of order two: its numerator does not vanish, and its
only vanishing denominator factor is $(t-1)^2$.
Both $u$ and $\lambda$ have poles of order two, with
respective leading terms $-h/20$ and $h/2$.
They are actual eigenvalue functions.
No change of finite free lattice can make $F$ regular
there: a matrix over the valuation ring has integral
eigenvalues, whereas these have negative valuation.
The finite localization explicitly removes $t=+1$
while retaining $t=-1$.
\end{proof}

The construction asserts neither a full-centralizer
identification nor all ordinary left morphisms.
Under $v=q^{1/2}$ with base $q=4$, the marked base
eigenvalues have weights $0,\tfrac12,1$ and different
absolute values. It therefore lies outside both the
common-radius and integral-mixed-weight hypotheses
of the marked-module gate.
It is an algebraic matrix family, not an assertion
of an ordinary geometric mixed object with integral
Deligne weights.
Its regularity at one place is also not a comparison
with VIII's prescribed global integral form or a
continuous choice along the positive real parameter
axis. The ordinary trace and marked trace remain
regular at $t=+1$ through the cancellations in
\eqref{eq:fractionalalgebraictraces}, even though the
actual operator has poles there.

\subsection{Regularity at every collision place}
The preceding exception is eliminated by an explicit
stronger regularity requirement. It does not require
integral weights or a common spectral radius.

\begin{finding}[\fid{F55}, Part 2: the all-places marked gate]
\label{f:allplacestrace}
Let $\mathcal L/\C(v)$ be a finite coefficient field
for an algebraic operator $F$ on the actual rank-$64$
tensor module.
Assume the regular-evaluation, base invertibility and
right-action commutation hypotheses of \fid{F54},
Part~3, and its marked-$X$ trace identity for all
sufficiently large extension degrees.
If $F$ is regular on the pulled-back actual tensor
lattice at every place of $\mathcal L$ above $v=1$,
where $X$ specializes to zero, this is impossible.
No weight or common-radius hypothesis is required.
\end{finding}

\begin{proof}
First let $f$ be any algebraic function over $\C(v)$
whose selected regular values satisfy
\[
 f_{b^n}=\sum_i m_i\alpha_i^n
 \quad\text{for all sufficiently large }n,\qquad
 b>1,\quad m_i\in\Z_{>0},\quad\alpha_i\ne0.
\]
Put $N=\sum_i m_i$.
By \fid{F54}, Part~1, there are $L>0$ and integers
$a_i$ with $\alpha_i^L=b^{a_i}$.
Let $P(v,Y)$ be the irreducible minimal polynomial of
$f$ over $\C(v)$, and put
\[
 p(t)=\sum_i m_i t^{a_i}.
\]
At the indices $n=Lr$, the selected trace values
are $p(b^r)$. The rational function $P(t^L,p(t))$
vanishes at infinitely many $t=b^r$, away from its
finitely many possible poles. It is identically zero.
The assignments
\[
 v\longmapsto t^L,\qquad f\longmapsto p(t)
\]
therefore give a field embedding
$\C(v,f)\hookrightarrow\C(t)$.
Restrict the $t=1$ place to this field.
Since $t^L-1$ has a simple zero at $1$, it is a place
over $v=1$ with value group $\Z$.
The function $f$ is regular there and has residue
$p(1)=N$.
This proves existence of one such place, not that
every place has this residue.

This trace-field place extends to the entire finite
coefficient field. Indeed, let $\mathcal O$ be its
discrete valuation ring and $\mathcal B$ the integral
closure of $\mathcal O$ in $\mathcal L$.
The extension is separable in characteristic zero.
The finite-normalization theorem makes
$\mathcal B$ finite over $\mathcal O$
(Stacks Lemma~10.161.8,
\href{https://stacks.math.columbia.edu/tag/032L}{tag~032L}).
Lying over gives a prime $\mathfrak q$ above the
maximal ideal of $\mathcal O$
(Stacks Lemma~10.36.17,
\href{https://stacks.math.columbia.edu/tag/00GQ}{tag~00GQ}).
The localization $\mathcal B_{\mathfrak q}$ is a
one-dimensional Noetherian normal local domain,
hence a discrete valuation ring
(Stacks Lemma~10.119.7,
\href{https://stacks.math.columbia.edu/tag/00PD}{tag~00PD}).
Here dimension one follows also because the generic
fiber of this finite integral domain extension is
the field $\mathcal L$.
Its residue field is finite over $\C$, thus equals
$\C$. The extended place retains the residue $f=N$.

Now apply this to
\[
 f(v)=\Delta(b)^{-1}\Tr_M(X(v)F(v)).
\]
At $b>1$, the marked sector has dimension eight.
Commutation with the actual right action makes it
invariant under the invertible $F_b$.
The marked comparison identifies the selected values
of $f$ with the genuine power traces of that
eight-dimensional restriction.
The construction just proved gives a place of its
trace field with residue $8$ and extends it to
$\mathcal L$.

By the all-places hypothesis, $F$ is regular on the
actual tensor lattice at this place and $X$ specializes
to zero. Therefore $\Tr(XF)$, and hence $f$, has
residue zero. This contradicts residue $8$.
\end{proof}

The same argument applies if, at every relevant place,
$X$ specializes to a nilpotent operator commuting
with the regular specialization of $F$:
then $XF$ is nilpotent and has trace zero.
It also applies to any fixed nonzero marked-sector
rank, replacing eight by that rank.
Regularity at every place is additional data;
it is not forced by the mere existence of an
algebraic coefficient scheme.
Part~1 violates exactly this condition at $t=+1$
while satisfying the chosen negative-place
regularity, the actual right action and both
trace laws. Its normalized marked function has
$g(-1)=0$ and $g(1)=8$.

% Separate paragraph for Part III.
\paragraph{\fid{F55}, Parts 1--2, in Grothendieck's language.}
The constructive model is a finite coefficient cover
$v=t^2$, followed by an explicitly specified finite
localization, carrying the actual pulled-back
coherent algebra, tensor module and right-module
automorphism.
The affine module/sheaf dictionary is EGA~I,
Theorem~1.3.7, p.~86, Corollary~1.3.8, p.~87,
and Theorem~1.4.1, pp.~90--92.
The representation and commuting-morphism language
is SGA~3, I, 4.7.2--3, with cote~90, scan~17 as the
audited framework anchor \cite{EGA,SGA3,Archive}.
The all-places argument uses the effective torus
character sum after $v=t^L$ from \fid{F54};
its positive multiplicities belong to the effective
cone in the representation Grothendieck group,
with SGA~6, V, \S\S2--3 and cote~133, scan~18 as
the framework anchors \cite{SGA6,Archive}.
The precise finite-normalization, lying-over and
discrete-valuation facts are the three primary
Stacks lemmas cited in the proof; the general
EGA dictionary is not used as their specific locator.
The two places above $v=1$ and the operator pole
between their regularity claims are explicit
algebraic data.
A continuous scalar branch along the positive real
locus is a different, additional comparison,
treated separately in Part~3.
Neither the coherent model nor these trace identities
choose the designated canonical basis, global
integral form, or target positivity inequalities.

